% Part: lambda-calculus
% Chapter: introduction
% Section: composition

\documentclass[../../../include/open-logic-section]{subfiles}

\begin{document}

\olfileid[hi]{lam}{int}{com}
\olsection{\usetoken{S}{lambda definable} फलन संयोजन के अधीन संवृत हैं}

\begin{lem}
!!{lambda definable} फलन संयोजन के अधीन संवृत हैं।
\end{lem}

\begin{proof}
मान लें $f$ को $h$, $g_0,$ \dots,~$g_{k-1}$ से संयोजन द्वारा परिभाषित किया
गया है। यदि $h$, $g_0$, \dots,~$g_{k-1}$ क्रमशः $H$, $G_0$,
\dots,~$G_{k-1}$ द्वारा !!{lambda defined} हैं, तो हमें ऐसा पद~$F$ खोजना है
जो~$f$ को !!{lambda define} करे। किंतु हम~$F$ को सीधे इस प्रकार परिभाषित
कर सकते हैं:
\[
F(x_0, \dots, x_{l-1}) = H(G_0(x_0, \dots, x_{l-1}),
\dots, G_{k-1}(x_0, \dots, x_{l-1})).
\]
दूसरे शब्दों में, लैम्ब्डा कलन की भाषा संयोजन को निरूपित करने के लिए अत्यंत
उपयुक्त है।
\end{proof}

\end{document}
